%%=============================================================================
%% De inhoud van je projectvoorstel
%%=============================================================================
%%
%% Zeven rubrieken, samen hoogstens twee bladzijden. Houd het kort: per
%% rubriek volstaan één tot drie zinnen. Een voorstel is geen rapport - het
%% moet je begeleider in vijf minuten laten beslissen of dit een geschikt
%% graduaatsproject is.
%%
%% Stem je voorstel af met je mentor VOORDAT je het indient. Je mentor
%% bevestigt dat dit is wat het bedrijf bedoelt en dat de tien werkdagen
%% haalbaar zijn.

%---------- 1. Context en probleemstelling ------------------------------------

\section{Context en probleemstelling}%
\label{vst:context}

% Waar loop je stage, wat doet dat team, en welk concreet probleem of welke
% concrete behoefte is de aanleiding? Noem de mensen voor wie het een probleem
% is. Vermijd algemeenheden als "bedrijven hebben nood aan".

Normaal komt de opdracht voor het graduaatsproject van het stagebedrijf; omdat
ik stage loop bij HOGENT, gaf mijn lector van het vak Graduaatsproef mij deze
losstaande opdracht. Ik migreer een afgebakend deel van PyPortfolioOpt naar
C\#/.NET; de Pythonbibliotheek berekent rendement en risico voor portefeuilles
en vormt de referentie~\autocite{PyPortfolioOptDocs}.

%---------- 2. Doelstelling ---------------------------------------------------

\section{Doelstelling}%
\label{vst:doelstelling}

% Wat bestaat er op het einde van dit project dat er vandaag niet is? Formuleer
% het als een resultaat waarvan iemand kan vaststellen dat het er is, en
% vermeld waaraan je zal merken dat het geslaagd is.

Ik lever een zelfstandige .NET-toepassing die koersgegevens uit CSV-bestanden
leest uit historische of gesimuleerde data. Ze berekent historisch rendement
en risico met steekproef- of shrinkagecovariantie. Optimalisaties richten zich
op minimumvolatiliteit, maximale Sharpe-ratio of een doelrendement. De app
rapporteert prestaties, normaliseert gewichten en ondersteunt discrete
allocatie; tests vergelijken de resultaten met Python.

%---------- 3. Technologieën --------------------------------------------------

\section{Technologieën}%
\label{vst:technologieen}

% Met welke talen, frameworks, databanken en diensten ga je werken? Welke
% liggen vast door de omgeving van het bedrijf, en welke kies je zelf? Vermeld
% ook wat nieuw voor je is: technologische verbreding wordt gewaardeerd.
%
% Vermeld hier ook of je AI denkt in te zetten, en waarvoor. Dat mag en wordt
% zelfs verwacht; in je rapport verantwoord je later wat je precies gedaan hebt.
%
% Verwerkt je toepassing persoonsgegevens, of zit er AI in wat je bouwt? Zeg
% dat dan nu al in één zin: dan weet je begeleider meteen dat de AVG of de
% AI Act meespeelt, en kan dat mee bekeken worden bij de goedkeuring.

Ik gebruik C\# met .NET 10, MathNet.Numerics, xUnit en een optimalisatiesolver;
de architectuur volgt Domain-Driven Design~\autocite{MathNetDocs,DotNetDocs}.
Python dient alleen als testreferentie, zodat de app zonder Python draait. AI
helpt bij analyse, ontwerp, code en documentatie; ik controleer de uitkomsten
en documenteer minstens acht interacties en twee gecorrigeerde fouten in
AI-USAGE.md. De app verwerkt geen persoonsgegevens, voert geen transacties uit
en geeft geen persoonlijk beleggingsadvies.

%---------- 4. Afbakening -----------------------------------------------------

\section{Afbakening}%
\label{vst:afbakening}

% DE BELANGRIJKSTE RUBRIEK. Je stage en je graduaatsproject lopen bij hetzelfde
% bedrijf; dit is waar je ze uit elkaar houdt.
%
% Beschrijf:
%   - wat je dagelijkse stagewerk is;
%   - wat het afgebakende geheel is dat JIJ zelfstandig analyseert, ontwerpt
%     en oplevert - dat is je graduaatsproject;
%   - wat er uitdrukkelijk buiten valt.
%
% Toets voor jezelf: welke analyse en welk ontwerp maak jij zelf? Kan je daar
% niets over schrijven, dan heb je een taak en geen project.

Mijn dagelijkse stageopdrachten bij HOGENT vallen buiten dit project. Ik
analyseer en ontwerp zelfstandig de migratie van de opgegeven functies en lever
de .NET-app, tests en vergelijking met Python op; een volledig financieel
platform, live beurskoppelingen, transacties en persoonlijk advies vallen
buiten scope.

%---------- 5. Verwachte meerwaarde -------------------------------------------

\section{Verwachte meerwaarde}%
\label{vst:meerwaarde}

% Wat heeft het bedrijf hieraan? Wat blijft er van jouw werk over wanneer je
% stage afgelopen is? Denk aan tijdwinst, minder fouten, nieuwe mogelijkheden,
% lagere kosten, of een basis waarop verder gebouwd kan worden.

HOGENT krijgt een demonstrator, broncode en tests als basis voor de
portefeuilleberekeningen in .NET. De vergelijking toont waar de resultaten
overeenkomen en verklaart eventuele numerieke verschillen.

%---------- 6. Planning en werkbelasting --------------------------------------

\section{Planning en werkbelasting}%
\label{vst:planning}

% Het graduaatsproject is ongeveer tien werkdagen, samen ongeveer 100 uren.
% Verdeel die over de fasen van je project en geef per fase aan wat je oplevert.
% Vermeld ook of je die dagen gefocust opneemt of gespreid over de stageperiode,
% en dat je dit met je mentor afgesproken hebt.

\begin{table}[htbp]
  \centering
  % tabularx met X-kolommen: de tabel is precies zo breed als de kolom (\linewidth) en
  % lange tekst breekt af. Een gewone tabular met l-kolommen werd in de tweekolomsopmaak
  % breder dan de kolom en liep in de rechterkolom (Overfull \hbox in het logboek).
  \begin{tabularx}{\linewidth}{@{}>{\raggedright\arraybackslash}X>{\raggedright\arraybackslash}X r@{}}
    \toprule
    \textbf{Fase} & \textbf{Wat lever je op} & \textbf{Uren} \\
    \midrule
    Analyse en vereisten      & scope en gedrag Python-referentie       & 8   \\
    Ontwerp                   & architectuur en solverkeuze             & 10  \\
    Realisatie                & .NET-berekeningen en optimalisatie      & 32  \\
    Testen en validatie       & tests en vergelijking met Python        & 12  \\
    Documentatie en demo      & AI-log, verslag en presentatie          & 8   \\
    \midrule
    \textbf{Totaal}           &                              & \textbf{70} \\
    \bottomrule
  \end{tabularx}
  \caption[Planning en werkbelasting]{\label{tab:vst-planning}%
    Voorlopige planning van 70 uur (raming opdrachtfiche: 60--75 uur), nog af te stemmen met de mentor.}
\end{table}

Ik stel voor de uren te spreiden over de stageperiode, in overleg met mijn
mentor. Ik plan minstens drie herkenbare Git-checkpoints en een presentatie
van 12--15 slides met een demo van 15 minuten en 5 minuten vragen.

%---------- 7. Risico's -------------------------------------------------------

\section{Risico's}%
\label{vst:risicos}

% Wat kan dit project doen mislukken, en wat doe je dan? Denk aan: toegang tot
% gegevens of systemen die je nodig hebt, afhankelijkheid van een collega of
% een derde partij, technologie die je nog moet leren, of drukte op de
% werkvloer waardoor je dagen niet vrijkomen.
%
% Twee tot drie risico's met telkens één zin over wat je terugvalplan is.

Een optimalisatiesolver kan extra integratiewerk vragen; daarom kies ik die
vroeg en valideer ik eerst minimumvolatiliteit en Sharpe-ratio. Numerieke
verschillen of ontbrekende historische data vang ik op met tolerantiegrenzen en
gesimuleerde testgegevens. De voorgestelde 70 uur en spreiding over de
stageperiode stem ik nog met mijn mentor af.
